\section{Introduction}
\label{sec:introduction}

Learning from model-generated answers can reduce the need for human
demonstrations. STaR and reinforced self-training alternate generation,
selection, and parameter updates \citep{zelikman2022star,singh2023restem}.
This feedback is a component of recursive self-improvement: the updated
learner shapes the candidates available for subsequent training.
Accepted answers become training targets, so verifier mistakes can affect
both the next update and the candidate distributions encountered later.

On a common candidate distribution, true- and false-acceptance rates determine
yield and precision, but not which errors are accepted. Their gradients can
differ despite equal counts. Here, \emph{influence} denotes an example's
contribution to the training gradient, including its magnitude and direction.

We ask what matching these statistics leaves undetermined about learning.
For a scalar logistic learner with group magnitudes $1$ and $L$, two rules
share all four initial summaries yet induce opposite accuracy changes when
$L>3$ in our illustrative parameter setting. We also prove long-run
separation under resampled population gradient flow and a sharp uniform
threshold for a specified group-targeted audit. These construction-specific
results imply neither universal separation nor equivalent AdamW dynamics.

For language models, we jointly construct final training subsets to control
prompt coverage, retention, and training volume. Both branches share the
candidate pool, initialization, prompts, correct responses, error counts,
and, for graphs, supervised-token counts per prompt. Only error selection
changes; coincident random and structured choices are retained. Graph errors
target edge-valid nonshortest routes; arithmetic errors target ignored
parentheses. Exact judges identify these signatures, not gradient influence
(Section~\ref{sec:tasks}).

\paragraph{Our contributions.}
\begin{itemize}
\item \textbf{Separation and auditing.} We establish population separation
at matched acceptance rates and a sharp corrective threshold for a
specified group-targeted audit.
\item \textbf{From gradients to updates.} Matched-subset identities isolate
accepted-error gradients; smooth-loss bounds connect them to parameter
updates while distinguishing reference loss from task risk
(Propositions~\ref{prop:finite-gap}--\ref{prop:remainder}; proofs in
Appendix~\ref{app:proof}).
\item \textbf{Controlled empirical comparisons.} Five-block graph studies
across three language models compare four-round audited and unaudited
self-training with isolated one-step updates. Qwen3-1.7B exhibits differing
four-round error composition; primary one-step accuracy contrasts are small.
\end{itemize}
